\documentclass[11pt]{article}
\usepackage[a4paper,margin=2.35cm]{geometry}
\usepackage{fontspec}
\usepackage{amsmath,amssymb}
\setmainfont[ItalicFont=Noto Serif Italic]{Malgun Gothic}
\setsansfont{Malgun Gothic}
\setlength{\parindent}{1.2em}
\setlength{\parskip}{0.45em}
\setlength{\emergencystretch}{3em}
\begin{document}

2. \textbf{유한성 판정기준의 증명.} 이 조건이 필요함은 명백하다. 실제로 $\mathfrak S$가 유한 생성 정역이면, $\mathfrak S$ 자체를 판정기준에 나타나는 유한 생성 부분환 $\mathfrak R$로 잡을 수 있다.

역으로 계수체의 표수가 $p$일 때에는 \textbf{근체(Wurzelkörper) $P^{1/p}$가 계수체 $P$의 유한 확대라는 가정}\footnote{원논문 28쪽 각주 2 참조.} 아래에서 이 조건이 충분하다. 표수가 0일 때에는 아무 제한 가정도 필요하지 않다. 이를 증명하려면 $\mathfrak S$가 $\mathfrak R$의 어떤 부분환에 관한 \textbf{유한 가군 생성계(Modulbasis)}를 가짐을 보이면 된다.

$\mathfrak K$는 $\mathfrak R$의 분수체\footnote{원논문 28쪽 각주 2 참조.}를, $\mathfrak L$은 $\mathfrak S$의 분수체를 뜻한다고 하자. $\mathfrak R$과 $\mathfrak S$는 영인자가 없는 환이므로 이 두 분수체가 존재하며, $P<\mathfrak K<\mathfrak L<P(x_1,\ldots,x_n)$이다. 따라서 1의 따름정리에 의해 $\mathfrak L$은 $\mathfrak K$의 \textbf{유한 대수적 확대체}다. 또한 가정에 의해 $\mathfrak S$의 모든 원소는 $\mathfrak R$에 대해 정수적이므로, $\mathfrak S$는 $\mathfrak L$에서 $\mathfrak R$에 대해 정수적인 모든 원소로 이루어진 환 $\mathfrak S'$의 부분환이다. 그러나 $\mathfrak R$이 $\mathfrak K$ 안에서 정수적으로 닫혀 있다고 가정하지 않았으므로, 유한 확대에서 약수사슬정리(Teilerkettensatz)가 보존된다는 사실을 직접 적용할 수는 없다. 그 대신 먼저, 역시 유한 생성인 $\mathfrak R$의 부분환 $\mathfrak T$로 넘어가야 하며, 이 $\mathfrak T$는 그러한 정수적 닫힘 조건을 만족해야 한다.

\end{document}
